% source: https://www.overleaf.com/learn/latex/Code_listing

\documentclass{article}
\usepackage[utf8]{inputenc}

\usepackage{listings}
\usepackage{xcolor}

%New colors defined below
\definecolor{codegreen}{rgb}{0,0.6,0}
\definecolor{codegray}{rgb}{0.5,0.5,0.5}
\definecolor{codepurple}{rgb}{0.58,0,0.82}
\definecolor{backcolour}{rgb}{0.95,0.95,0.92}

%Code listing style named "mystyle"
\lstdefinestyle{mystyle}{
  backgroundcolor=\color{backcolour},   commentstyle=\color{codegreen},
  keywordstyle=\color{magenta},
  numberstyle=\tiny\color{codegray},
  stringstyle=\color{codepurple},
  basicstyle=\ttfamily\footnotesize,
  breakatwhitespace=false,
  breaklines=true,
  captionpos=b,
  keepspaces=true,
  numbers=left,
  numbersep=5pt,
  showspaces=false,
  showstringspaces=false,
  showtabs=false,
  tabsize=2
}

%"mystyle" code listing set
\lstset{style=mystyle}


\begin{document}
\maketitle
\thispagestyle{empty}
\newpage

\pagenumbering{roman}
\tableofcontents
\newpage

\pagestyle{plain}
\pagenumbering{arabic}

\section{Introduction}
\label{sec:orge6a1037}
This example shows a lot of what you can achieve with \texttt{doc-org}.

However, \texttt{doc-org} uses latex, emacs and org-mode, so the sky is the limit and
there will never be a \emph{"full"} example.

\subsection{header}
\label{sec:org4bdb7bb}
The \texttt{header.tex} file of this example is a good starting point for your
document, because it supports a lot of features.
Of course, reading all the other examples related to each single feature is
recommended to understand what's the role of every line of the header.
\section{some code}
\label{sec:orgb279d2c}
If you specify the programming language, code syntax is highlighted.
\begin{minted}[]{c}
int main()
{
    printf("hello world\n");
}
\end{minted}

\begin{minted}[]{python}
# This program prints Hello, world!
print('Hello, world!')
\end{minted}

\section{Some math}
\label{sec:orgcbab23d}
You can insert latex equation, like equation \ref{eq:org5da5461}.
\begin{equation}
\label{eq:org5da5461}
\phi = \frac{2\pi fD}{c}
\end{equation}
\subsection{Other latex equations}
\label{sec:orge394543}
Equation \ref{eq:orgcfc0044} reference.

\begin{equation}
\label{eq:orgcfc0044}
D = \frac{c\phi}{2\pi f}
\end{equation}

\subsubsection{A more complicated equation}
\label{sec:org756a47b}

\begin{equation}
\Delta TOF_{est} = \frac{k_T TOF}{1+k_T } - 0.5 \frac{\mu_A' - \mu_T'}{1+k_T}.
\end{equation}

\section{lists}
\label{sec:orgbe76d66}
\subsection{numerated}
\label{sec:org0b1ecaf}
\begin{enumerate}
\item one
\item two
\begin{enumerate}
\item sub-two
\begin{enumerate}
\item sub-sub-two
\end{enumerate}
\end{enumerate}
\item \label{org46fea8c}three
\end{enumerate}

Here we refer to item \ref{org46fea8c}.

\subsection{unnumerated}
\label{sec:org9fd951f}
\begin{itemize}
\item One
\begin{itemize}
\item Two
\begin{itemize}
\item Three
\begin{itemize}
\item Four
\end{itemize}
\item Five
\end{itemize}
\end{itemize}
\end{itemize}

\subsubsection{trivia}
\label{sec:org01cced7}
Doc-org could mean a lot of things, such as:
\begin{itemize}
\item document - organized
\item docker - org mode
\item docile - organist
\end{itemize}
A team of 42 linguists is currently searching for the answer.

\subsection{Enumeration example}
\label{sec:orga4f0120}
\begin{description}
\item[{element1}] description for element one.
\item[{element2}] Very long description for element two, in order to show a
multiline description.
\end{description}
\section{tables}
\label{sec:org7fa93c2}

look how cute is this simple table:

\begin{center}
\begin{tabular}{rr}
one & two\\
\hline
1 & 2\\
3 & 4\\
\end{tabular}

\end{center}

\begin{table}[htbp]
\centering
\begin{tabular}{lrrrrr}
Anchor & MAE & Correlation & MMA & MAPE & RMSE\\
\hline
a1 & 6.871 & -0.372 & 1.164 & 0.157 & 8.515\\
a2 & 5.275 & 0.329 & 1.127 & 0.117 & 6.332\\
a3 & 5.406 & 0.672 & 1.117 & 0.114 & 6.457\\
\end{tabular}
\caption{\label{tab:org9ab42ba}
You can add captions to your tables}

\end{table}

\section{Images}
\label{sec:org9940e54}
You can insert image, for example see image \ref{fig:orgb8e5d04}.

\begin{figure}[htbp]
\centering
\includegraphics[width=.9\linewidth]{./img/example_image.png}
\caption{\label{fig:orgb8e5d04}
Example image.}
\end{figure}

\subsection{image matrix}
\label{sec:org58d036b}
In image \ref{fig:image_matrix} you can find an example of how to combine different images
into one in latex.

\begin{figure}[!tbp]
  \centering
  \subfloat[Case 1.]{\includegraphics[width=0.55\textwidth]{./img/example_image.png}\label{fig:f1}}
  \subfloat[Case 2.]{\includegraphics[width=0.55\textwidth]{./img/example_image.png}\label{fig:f2}}
  \\
  \subfloat[Case 3.]{\includegraphics[width=0.55\textwidth]{./img/example_image.png}\label{fig:f3}}
  \subfloat[Case 4.]{\includegraphics[width=0.55\textwidth]{./img/example_image.png}\label{fig:f4}}
  \\
  \subfloat[Case 5.]{\includegraphics[width=0.55\textwidth]{./img/example_image.png}\label{fig:f5}}
  \subfloat[Case 6.]{\includegraphics[width=0.55\textwidth]{./img/example_image.png}\label{fig:f6}}
  \caption{\label{fig:image_matrix}Combine more pictures into one.}
\end{figure}

\end{document}
